\documentclass[11pt]{letter}

\usepackage[letterpaper,margin=1in]{geometry}
\usepackage[hidelinks]{hyperref}
\usepackage[T1]{fontenc}
\usepackage{lmodern}

\longindentation=0pt
\setlength{\parindent}{0pt}
\setlength{\parskip}{0.6em}

\signature{Austin Jiang}

\begin{document}

{\Large \textbf{Austin Jiang}}\\
\href{mailto:a68jiang@uwaterloo.ca}{a68jiang@uwaterloo.ca} \textbar\ 
\href{https://github.com/AustinBoyuJiang}{github.com/AustinBoyuJiang} \textbar\ 
\href{https://www.linkedin.com/in/austin-boyu-jiang}{linkedin.com/in/austin-boyu-jiang}

\vspace{0.8em}

Dear hiring manager,

I am a Computer Science student at the University of Waterloo applying for the Software Engineering Intern role. My interests lie in performance engineering, concurrent systems, and machine learning systems.

I am currently an undergraduate research assistant in the Multicore Lab at the University of Waterloo, where I work on benchmarking and extending the \textbf{Verlib (PPoPP’24)} system. My work includes migrating the build to the CUDA toolchain to resolve host-device compatibility issues, developing an adapter supporting versioned and non-versioned data structures, and benchmarking range queries at \textbf{10M-key scale under 100+ threads}. I also implemented B+ Tree and ART data structures using the GPU Range Query Dictionary library for performance comparison.

Previously at Wolfram Research, I developed a multi-threaded engine managing over 800K data points across 8 CPU cores and designed sparse frontier updates and dense scans, achieving a \textbf{127\% speedup} through improved load balancing. I also built a benchmarking harness measuring throughput. My early work has also produced a first-author research paper presented at the Wolfram Technology Conference 2025.

Recently, I was shortlisted for \textbf{Mistral AI’s core research internship}. Although the process was paused due to headcount, the technical evaluation confirmed my understanding of KV cache optimization, RoPE implementations, and memory-efficient Transformers. In addition, my TreeHacks project involved deploying Hunyuan WorldPlay 1.5 across 4× H200 GPUs using tensor parallelism.

My competitive programming background, including being \textbf{two-time Canadian Computing Olympiad top 10 nationally} and USA Computing Olympuad highest division, showcases my strong problem-solving ability in algorithms and data structures.

I would love the opportunity to contribute to your team.

Sincerely,\\
Austin Jiang

\end{document}